\chapter{Proofs for the Base-Direction Classifications}
\label{app:base-proofs}

This appendix proves the rank and quotient statements and the intrinsic defect classification of
\Cref{ch:pullback-geometry}, and collects four witnesses that mark the boundaries of the
classification. The notation is that of \Cref{sec:pb-rank-quotient} and
\Cref{sec:pb-intrinsic-defect-classification}.

\section{Rank and quotient geometry}
\label{sec:app-base-rank-quotient}

\noindent\emph{Proof of \Cref{thm:pb-pullback-rank-quotient}.}
At each \(c\) the jet \(D^\omega s\) is a linear map \(T_c\mathcal U\to V_{s(c)}E\) and
\(h_s^\omega=(D^\omega s)^*g^F\). The radical of the pullback of a positive-semidefinite form by a
linear map is the preimage of its radical, which is the pointwise algebra of
\Cref{thm:hist-retained-fisher-quotient}; since \(g^F\) is positive definite, the preimage is
\(\ker D^\omega s\), and the rank statement follows. The constant-rank statement is the one for
morphisms of smooth vector bundles over a fixed base, not the constant-rank theorem for maps of
manifolds: \(D^\omega s\) is a morphism \(T\mathcal U\to s^*VE\) over the identity, and a morphism of
locally constant rank has smooth kernel and image subbundles by local adaptation of frames. Replacing
an argument of \eqref{eq:pb-pullback-quotient-metric} or \eqref{eq:pb-pullback-quotient-amari} by a
representative that differs by an element of \(K_s^\omega\) leaves its vertical jet unchanged, so the
quotient tensors are well defined. The factored map is injective onto the image of \(D^\omega s\) and
preserves the displayed tensors, so the quotient metric is positive definite. \hfill$\square$

\noindent\emph{Proof of \Cref{prop:pb-quotient-manifold-requirements}.}
(i) If \(q_K\) exists with \(\ker Tq_K=K_s^\omega\), the distribution is tangent to the fibers of
\(q_K\), which are submanifolds, so brackets of its sections are again its sections; since
\(K_s^\omega=\ker D^\omega s\), this is \eqref{eq:pb-null-involutivity}. For the witness, take the
torus \(\mathbb T^2=\R^2/\mathbb Z^2\), the trivial principal translation bundle with fiber
\(\{\mathcal N(\mu,1)\}\), the constant section \(\mu=0\), and the connection form
\(A=\mathrm dy-\lambda\mathrm dx\) with \(\lambda\) irrational. Then \(D^\omega s=A\) has constant
rank one, its kernel is spanned by \(\partial_x+\lambda\partial_y\) and is involutive, and every leaf
is the image of a line of irrational slope, which is dense in \(\mathbb T^2\) and not all of it. The
fibers of a continuous map into a Hausdorff space are closed, and a dense proper leaf is not, so no
surjective submersion onto a Hausdorff manifold has the leaves as its fibers.

(ii) If \(\tau\) is a tensor on \(\overline{\mathcal U}\) with \(q_K^*\tau=h_s^\omega\), every
\(Z\in\Gamma(K_s^\omega)\) is \(q_K\)-related to the zero field, so its flow preserves the fibers and
\(\mathcal L_Zq_K^*\tau=0\); the same holds for the cubic tensor. For the witness, take \(\R^2\), the
trivial translation bundle with fiber \(\{\mathcal N(\mu,1)\}\), the zero section, and the connection
form \(A=\varphi(y)\mathrm dx\) with \(\varphi\) smooth, positive, and nonconstant. Then
\(D^\omega s=\varphi(y)\mathrm dx\), the null leaves are the vertical lines, \(q_K(x,y)=x\) is a smooth
Hausdorff quotient, and \(h_s^\omega=\varphi(y)^2\mathrm dx^2\) has
\(\mathcal L_{\partial_y}h_s^\omega=2\varphi\varphi'\mathrm dx^2\neq0\) where \(\varphi'\neq0\); no
tensor \(\tau(x)\mathrm dx^2\) on \(\R\) pulls back to it.

(iii) For \(Z\in\Gamma(K_s^\omega)\),
\((\mathcal L_Zh_s^\omega)(X,Y)=Z(h_s^\omega(X,Y))-h_s^\omega([Z,X],Y)-h_s^\omega(X,[Z,Y])\). Metric
compatibility gives
\(Z\bigl(g^F(D^\omega sX,D^\omega sY)\bigr)=g^F(\nabla_Z(D^\omega sX),D^\omega sY)+g^F(D^\omega sX,\nabla_Z(D^\omega sY))\),
and subtracting the bracket terms gives \eqref{eq:pb-basicness-criterion}. If
\(\nabla'=\nabla+B\) is another \(g^F\)-compatible connection, \(B_Z\) is \(g^F\)-skew-adjoint, so the
changes of the two summands cancel; the left side is the ordinary Lie derivative and does not depend on
\(\nabla\). \hfill$\square$

\section{Proofs for the intrinsic classification}
\label{sec:app-base-classification}

\noindent\emph{Proof of \Cref{thm:pb-defect-contractions}.}
The coarse jet is \(\mathsf p+\mathsf J\eta=\mathsf p_\perp+C\) with
\(C=\mathsf p_{\mathcal O}+\mathsf J\eta\) taking values in \(\mathcal O\). By orthogonality the left side
of \eqref{eq:pb-intrinsic-positivity-problem} is \(\mathsf p_\perp^*g_1\mathsf p_\perp+C^*g_1C\), so the
problem is \(C^*g_1C\preceq\mathsf H\). This forces \(\mathsf H\succeq0\) and \(Cx=0\) whenever
\(\mathsf H(x,x)=0\). For a positive-semidefinite form the diagonal null set is the radical, since
positivity of \(\mathsf H(x+ty,x+ty)\) for both signs of small \(t\) forces \(\mathsf H(x,y)=0\). Hence
\(\mathsf k([x])=Cx\) is well defined, and the inequality says exactly that \(\mathsf k\) is a
contraction; conversely every contraction satisfies it. Since
\(\mathsf k\operatorname{pr}_{\mathsf H}-\mathsf p_{\mathcal O}\) takes values in \(\mathcal O\), a linear
right inverse of \(\mathsf J\) on its image gives a lift, and two lifts differ by a map into
\(\ker\mathsf J\). If \(\mathsf H\succeq0\), \(\mathsf k=0\) proves existence. \hfill$\square$

\noindent\emph{Proof of \Cref{cor:pb-defect-markov-feasible}.}
Fiberwise, \Cref{thm:pb-pullback-fisher-defect} gives \(g_0-\Psi_*^*g_1\Psi_*\succeq0\), hence
\(\delta\succeq0\). Orthogonality gives
\(\mathsf p^*g_1\mathsf p=\mathsf p_\perp^*g_1\mathsf p_\perp+\mathsf p_{\mathcal O}^*g_1\mathsf p_{\mathcal O}\),
so \(\mathsf H=\mathsf D^*g_0\mathsf D-\mathsf p^*g_1\mathsf p+\mathsf p_{\mathcal O}^*g_1\mathsf p_{\mathcal O}\),
which is \eqref{eq:pb-defect-markov-budget}. The zero anomaly corresponds to \(C=\mathsf p_{\mathcal O}\),
and \(\mathsf p_{\mathcal O}^*g_1\mathsf p_{\mathcal O}\preceq\mathsf H\). For feasible \(C_0,C_1\) and
\(0\leq t\leq1\),
\((1-t)\lVert C_0x\rVert^2+t\lVert C_1x\rVert^2-\lVert(1-t)C_0x+tC_1x\rVert^2=t(1-t)\lVert(C_0-C_1)x\rVert^2\geq0\),
so the feasible set is convex, and the set of lifts is the preimage of a convex set under a linear map.
\hfill$\square$

\noindent\emph{Proof of \Cref{thm:pb-defect-partial-isometries}.}
In the quotient inner product,
\((\mathsf p+\mathsf J\eta)^*\Delta(\mathsf p+\mathsf J\eta)=(\mathsf p_\Delta+\mathsf J_\Delta\eta)^*(\mathsf p_\Delta+\mathsf J_\Delta\eta)\).
The component \((I-P_\Delta)\mathsf p_\Delta\) is unchanged by \(\mathsf J_\Delta\eta\) and is orthogonal
to \(\mathcal O_\Delta\), so removing its identical contribution from both sides reduces the cocycle to
\begin{equation}
(\mathsf w_\Delta+\mathsf J_\Delta\eta)^*(\mathsf w_\Delta+\mathsf J_\Delta\eta)=\mathsf w_\Delta^*\mathsf w_\Delta.
\label{eq:pb-defect-equal-gram}
\end{equation}
If \(\mathsf w_\Delta x=0\), equality of squared norms forces \((\mathsf w_\Delta+\mathsf J_\Delta\eta)x=0\),
so \(\mathsf T(\mathsf w_\Delta x)=(\mathsf w_\Delta+\mathsf J_\Delta\eta)x\) is well defined and linear,
and polarization shows that it is isometric. Conversely an isometric \(\mathsf T\) gives
\eqref{eq:pb-defect-equal-gram}. Every right side of \eqref{eq:pb-defect-isometry-parameterization} lies
in the image of \(\mathsf J_\Delta\), so it lifts with the stated kernel ambiguity. The choice
\(\mathsf T=\iota\) gives \(\mathsf J_\Delta\eta=0\), that is \(\mathsf J\eta\in\operatorname{rad}\Delta\),
which contains the isotropy case \(\mathsf J\eta=0\); \(\mathsf T=-\iota\) gives
\(-2\mathsf w_\Delta\); if \(\mathsf w_\Delta=0\) the only map is on the zero space; and \(\Delta=0\) makes
\(E_\Delta=0\). \hfill$\square$

\noindent\emph{Proof of \Cref{thm:pb-defect-signed-gram}.}
If \(o\in\mathcal O_1\) is \(\Delta\)-orthogonal to \(\mathcal O_1\), it is orthogonal to \(\mathcal O_0\)
as well, since \(\mathcal O_0\) is orthogonal to all of \(\mathcal O\); thus \(o\in\mathcal O_0\cap\mathcal O_1=0\),
which proves nondegeneracy of \(\Delta_1\) and uniqueness of \(\mathsf w\). Decompose
\(\mathsf J\eta=A_1+\mathsf n\) with \(A_1\in\mathcal O_1\) and \(\mathsf n\in\mathcal O_0\), and put
\(\mathsf k=\mathsf w+A_1\). Pairings within \(\mathcal O_0\) and between \(\mathcal O_1\) and
\(\mathcal O_0\) vanish, and each pairing of \(\mathsf p\) with \(\mathcal O_1\) may be replaced by the
pairing of \(\mathsf w\). Expanding the cocycle therefore gives exactly
\begin{equation}
\mathsf k^*\Delta_1\mathsf k=\mathsf G_{\mathsf n}.
\label{eq:pb-defect-signed-gram-factor}
\end{equation}
The terms with \(\mathsf l\) remain because the radical of \(\Delta|_{\mathcal O}\) need not be orthogonal
to \(\mathsf p\). The restriction of \(\mathsf G_{\mathsf n}\) to \(\mathcal V_{\mathsf n}\) is
nondegenerate, so a Gram factor restricts there to an injective isometry, whose image has signature
\((m_+,m_-)\) inside a space of signature \((n_+,n_-)\), and the inertia inequalities are necessary. Its
restriction to \(\mathcal K_{\mathsf n}\) is orthogonal to that image and totally isotropic. Conversely
these properties give the correct block on \(\mathcal V_{\mathsf n}\) and zero cross terms and block on
\(\mathcal K_{\mathsf n}\). The inertia inequalities suffice for \(\mathsf I_{\mathsf n}\): diagonalize
the two nondegenerate forms, rescale the eigenvectors to squared norms \(\pm1\), and map like-signed
basis vectors to each other; the spectral step maximizes a symmetric quadratic form on a unit sphere and
inducts on the orthogonal complement. Taking \(\mathsf Z=0\) gives existence. In a space of signature
\((r_+,r_-)\) the projections of a totally isotropic subspace onto the positive and the negative summands
are injective, since a vector with either projection zero and zero norm vanishes, so its dimension is at
most \(\min(r_+,r_-)\); applied to the nondegenerate complement of
\(\mathsf I_{\mathsf n}(\mathcal V_{\mathsf n})\), of signature \((n_+-m_+,n_--m_-)\), this bounds the rank
of \(\mathsf Z\). Lifting through \(\mathsf J\) completes the parameterization. In the example
\(\Delta(e_1,e_1)=0\), so \(\mathcal O_0=\mathcal O\), \(\mathcal O_1=0\), and \(\mathsf w=0\); then
\(\mathsf l_X(e_1)=X\Delta(e_2,e_1)=X\), so \(\mathsf G_{\mathsf n}(X,X)=-2\mathrm cX^2\), and
\(n_\pm=0\) forces \(\mathsf G_{\mathsf n}=0\); the cocycle difference is
\(2\mathrm cX^2\Delta(e_1,e_2)=2\mathrm cX^2\), which confirms \(\mathrm c=0\), and for \(\mathsf p=0\)
the form \(\mathsf l\) vanishes. \hfill$\square$

\noindent\emph{Proof of \Cref{prop:pb-defect-smooth-strata}.}
A nonzero minor of maximal size of a constant-rank smooth bundle map stays invertible locally; its
columns give a smooth image frame, and solving with it gives smooth preimages and kernel frames, so the
kernels, images, and quotients are smooth bundles. Positive bundle metrics are obtained by patching
local metrics with a partition of unity, orthogonal projections are smooth because their local matrices
invert positive Gram matrices, and restriction to the orthogonal complement of a kernel gives a smooth
isomorphism onto the image, hence a smooth right inverse of \(\mathsf J\) or \(\mathsf J_\Delta\). These
constructions also prove the converse: the maps
\(\mathsf k([x])=(\mathsf p_{\mathcal O}+\mathsf J\eta)x\) and
\(\mathsf T(\mathsf w_\Delta x)=(\mathsf w_\Delta+\mathsf J_\Delta\eta)x\) induced by a smooth solution are
smooth. In the signed case \(\mathsf I_{\mathsf n}\) and \(\mathsf Z\) are restrictions of the smooth
\(\mathsf k\) to the chosen smooth splitting, and local frames for a nondegenerate form of fixed signature
are obtained by retaining a basis with nonzero successive squared norms near a point, orthogonalizing,
and dividing by the square roots of their absolute norms; mapping like-signed vectors gives local smooth
embeddings under the inertia bounds, with no claim that independent local isometries glue. The
intrinsic statements follow because every construction uses only the tensor transformation laws of the
listed data, under which contraction, isometry, inertia, and total isotropy are preserved.
\hfill$\square$

\noindent\emph{Proof of \Cref{thm:pb-defect-connection-descent}.}
The difference of two principal connections is horizontal and equivariant, hence the tensorial form of a
unique \(\bar\eta\) on the target base; conversely adding such a form preserves the values of a
connection on fundamental vectors and its equivariance. Pullback by \(\mathcal P\) therefore gives
exactly \eqref{eq:pb-defect-affine-realization}. A pullback annihilates \(\ker Tf\) and satisfies the
equality on each fiber. Conversely define \(\bar\eta_y(Y)=\delta\eta_c(X)\) using any \(c\) over \(y\) and
any \(X\) with \(T_cfX=Y\); the two conditions make this independent of both choices and linear in
\(Y\). In local submersion coordinates there is a smooth local section \(\mathsf r\) of \(f\), and
\(\bar\eta=\mathsf r^*\delta\eta\) proves smoothness; surjectivity gives uniqueness. For a
diffeomorphism both conditions are automatic. A different lift changes \(\eta\) by a
\(\ker\mathsf J\)-valued one-form, which may restore the two conditions. \hfill$\square$

\section{Witnesses for the classification}
\label{sec:app-base-witnesses}

\propositionheading{Witnesses for the defect classification}{prop:pb-defect-witnesses}
(i) Action rank can obstruct a continuous lift. For trivial \(\mathrm{SO}(2)\) principal bundles over the
identity map of \(\R\), fine and coarse normal-location families on \(\R^2\) with covariances \(I\) and
\(2I\), the rotation-equivariant channel adding an independent \(\mathcal N(0,I)\) vector, zero source
connection, and related mean sections \(\mu(x)=(x^2,0)\), one has on the unit base vector
\begin{equation}
\mathsf D=(2x,0),\qquad \Psi_*=I,\qquad
g_0=I,\qquad g_1=\tfrac12 I,\qquad
\mathsf J_x(t)=(0,x^2t).
\label{eq:pb-defect-action-rank-example}
\end{equation}
The smooth anomaly \((0,x)\) lies in \(\operatorname{im}\mathsf J_x\) at every point and satisfies the
positivity problem, but its unique lift at \(x\ne0\) is \(\eta_x=1/x\), so no continuous connection
defect realizes it across zero, although both Fisher metrics are positive.
\par
(ii) A smooth defect need not have an ambient isometry extension. Let \(a(x)=e^{-1/x^2}\) for
\(x\ne0\) and \(a(0)=0\), \(\mu(x)=\int_0^xa(t)\mathrm dt\), scalar unit-variance normal-location
families at the first two levels, an identity first arrow, translation group \(\R\), identity base maps,
zero source connection, and an additive unit-normal noise second arrow, so that \(\Delta=\tfrac12\),
\(\mathsf p=a\), and \(\mathsf J=1\) relative to \(\mathrm dx\). The defect
\begin{equation}
\eta_x=
\begin{cases}
-2a(x),&x>0,\\
0,&x\leq0,
\end{cases}
\label{eq:pb-defect-flat-jet-example}
\end{equation}
is smooth and satisfies the cocycle and the first-arrow positivity with equality, while the isometry
\(\mathsf T\) of \eqref{eq:pb-defect-isometry-parameterization} is \(-1\) on the positive and \(+1\) on the
negative half-line and has no continuous extension as a field in \(\operatorname{End}(\R)\). On the
other hand, if continuous \(\mathsf W,\mathsf Y\) obey \(\mathsf Y^*\mathsf Y=\mathsf W^*\mathsf W\), the
induced isometry \(\mathsf T(\mathsf Wx)=\mathsf Yx\) is continuous on the image total space: in local
bundle coordinates, if \((c_n,z_n)\to(c_0,z_0)\) with \(z_n\in\operatorname{im}\mathsf W_{c_n}\) and
\(\mathsf W_{c_0}u_0=z_0\), then
\begin{equation}
\lVert\mathsf T_{c_n}z_n-\mathsf T_{c_0}z_0\rVert
\leq C_0\lVert z_n-\mathsf W_{c_n}u_0\rVert
     +\lVert\mathsf Y_{c_n}u_0-\mathsf Y_{c_0}u_0\rVert\longrightarrow0,
\label{eq:pb-defect-image-continuity}
\end{equation}
with \(C_0<\infty\) from the uniform local equivalence of continuous positive metrics with coordinate
norms.
\par
(iii) Smooth algebraic solutions can fail connection descent. For \(f(x,y)=x\) on \(\R^2\), trivial
translation principal bundles with the principal map \((x,y,g)\mapsto(x,g)\), identical scalar
unit-variance normal-location fibers, an identity first fiber map, zero mean sections, source connection
\(-\mathrm dy\), and target reference connection zero, one has \(\mathsf p=-\mathrm dy\) and
\(\eta_0=\mathrm dy\); the candidate \(\eta=2\mathrm dy\) satisfies positivity with equality and, with an
additive unit-normal noise second arrow, the cocycle, but violates
\eqref{eq:pb-intrinsic-kernel-constraint} and is realized by no target connection. For
\(\mathcal C=\R\times\{+,-\}\), \(f(x,\epsilon)=x\), translation-equivariant normal-location arrows with
variances \(1\to2\to3\), zero source and target reference connections, and related mean sections
\(\mu=x\), the choice \(\eta=0\) on the positive copy and \(\eta=-2\mathrm dx\) on the negative copy solves
both problems with constant ranks and \(\ker Tf=0\), and is realized by no target connection.
\par
(iv) Signed isotropic solutions occur outside the Markov domain. For unit-covariance normal-location
fibers on \(\R^2\), full translation action, a constant related mean section, zero source connection, an
identity first arrow, and a second arrow whose mean map has derivative
\(L_2=\operatorname{diag}(0,\sqrt2)\), one has \(\Delta=I-L_2^\top L_2=\operatorname{diag}(1,-1)\),
\(\mathsf p=0\), and \(\mathsf J=I\), and the constant target connection inducing the anomaly
\(X\mapsto(X,X)\) is a nonzero solution of the cocycle problem in the totally isotropic term of
\Cref{thm:pb-defect-signed-gram}, although the ambient radical of \(\Delta\) is zero; the second arrow is
equivariant for the translation homomorphism \(L_2\) and is not fiberwise Markov. \status{ESTABLISHED}

\noindent\emph{Proof.}
(i) The fundamental field of the rotation generator at the mean \((x^2,0)\) is \(t(0,x^2)\), the isotropic
covariances are rotation invariant, and \(g_1=\tfrac12I\) is the Fisher metric of the mean of
\(\mathcal N(\mu,2I)\). The anomaly is \(\mathsf J_x\) applied to \(1/x\) for \(x\ne0\) and is zero at
\(x=0\), where \(\mathsf J_0=0\), and
\(\tfrac12\lVert(2x,x)\rVert^2=\tfrac52x^2\leq4x^2=\lVert(2x,0)\rVert^2\). A continuous realizing defect
would equal \(1/x\) away from zero.

(ii) Every one-sided derivative of \(\eta\) at zero is a polynomial in \(1/x\) times \(e^{-1/x^2}\) and
tends to zero, so \(\eta\) is smooth. The coarse jet is \(\mathsf p+\eta=-a\) for \(x>0\) and \(a\) for
\(x<0\), so \((\mathsf p+\eta)^2=\mathsf p^2\) gives both problems with equality, since the first arrow is
the identity with unit Fisher metric. The isometry \(\mathsf T\) sends \(a\) to \(-a\) for \(x>0\) and to
\(a\) for \(x<0\). For the continuity estimate,
\(\mathsf T_{c_n}z_n-\mathsf T_{c_0}z_0=\mathsf T_{c_n}(z_n-\mathsf W_{c_n}u_0)+\mathsf Y_{c_n}u_0-\mathsf Y_{c_0}u_0\),
the isometry bounds the first term, and continuity of \(\mathsf W\) and \(\mathsf Y\) sends both terms to
zero.

(iii) With the sign convention of \Cref{def:pb-connection-difference}, the jet of a constant zero mean
section is the local connection form, so \(\mathsf p=-\mathrm dy\) and
\(\eta_0=\mathcal P^*0-(-\mathrm dy)=\mathrm dy\). The coarse algebraic jet of \(\eta=2\mathrm dy\) is
\(\mathrm dy\), with \((\mathrm dy)^2=\mathsf p^2\) and \(\Delta=1-\tfrac12=\tfrac12\), but it is nonzero on
\(\partial_y\in\ker Tf\), and \(\eta-\eta_0=\mathrm dy\) does not annihilate \(\partial_y\); the
translation action is injective, so there is no other lift. In the second example
\(\mathsf p=\mathrm dx\), \(\mathsf J=1\), \(\Delta=\tfrac12-\tfrac13=\tfrac16\), both branches satisfy
\((\mathsf p+\eta)^2=\mathsf p^2\) and \(\tfrac12(\mathsf p+\eta)^2\leq\mathsf p^2\), and
\eqref{eq:pb-defect-all-fiber-descent} fails because the two copies over one \(x\) carry \(0\) and \(-2\).

(iv) The coarse Fisher metric of a unit-covariance mean is \(I\), so the pulled-back metric is
\(L_2^\top L_2=\operatorname{diag}(0,2)\), and \((X,X)\Delta(X,X)^\top=X^2-X^2=0\). Here
\(\mathcal O=\R^2\) with \(\Delta_1\) of signature \((1,1)\), \(\mathsf w=0\), and
\(\mathsf G_{\mathsf n}=0\), so the solution is a totally isotropic \(\mathsf Z\). A fiberwise Markov
arrow has \(\Delta\succeq0\) by \Cref{thm:pb-pullback-fisher-defect}, which \(\operatorname{diag}(1,-1)\)
violates. All four witnesses use proper normal laws and normalized channels, and a one-point observation
of likelihood one may be appended with evidence one. \hfill$\square$
